\section*{Capítulo \ref{ch:b2:two-wave-interference} --- Interferencias de dos ondas}

\begin{solution}{exo:b2:two-wave-interference:1}
$i = \lambda D/a = \qty{2.4}{mm}$; $x/i = 4.2$; décima franja oscura en $9.5i = \qty{22.5}{mm}$;
con $a = \qty{2}{mm}$: \qty{0.47}{mm}; en agua, $\lambda/n$: \qty{1.8}{mm}.
\end{solution}

\begin{solution}{exo:b2:two-wave-interference:2}
$1.25I_0 \pm I_0$: $V = 0.8$. $1.01I_0 \pm 0.2I_0$: $V = 0.2$ --- el término de interferencia
va como el cociente de \emph{amplitudes}, $1/10$, no como el de intensidades. Con el filtro:
$V = 2\sqrt{0.5}/1.5 = 0.94$.
\end{solution}

\begin{solution}{exo:b2:two-wave-interference:3}
$2ne = \qty{665}{nm}$: brillante en $665/(p + \tfrac12)$: \qty{443}{nm} (azul violáceo);
oscuro en $665/p$: \qty{665}{nm} (rojo): una película azul. $e = \qty{1}{\micro m}$: $2ne = \qty{2660}{nm}$,
brillante en $760$, $591$, $484$ y \qty{409}{nm} --- cuatro colores a la vez, un
blanco rosado deslavado.
\end{solution}

\begin{solution}{exo:b2:two-wave-interference:4}
$r_p = \sqrt{p\lambda R}$: $1.09$, $1.53$, \qty{1.88}{mm}; $p = r^2/\lambda R = 85$ anillos dentro de
\qty{1}{cm}; el centro es oscuro: espesor nulo y una sola inversión de signo. Con agua:
radios divididos por $\sqrt n$, $98$ anillos.
\end{solution}

\begin{solution}{exo:b2:two-wave-interference:5}
(a) $i = 0.8$ y \qty{1.4}{mm}; la segunda roja, \qty{2.8}{mm}, y la tercera violeta, \qty{2.4}{mm}.
(b) El orden $p$ del rojo, en $1.4p$ mm, se encuentra con el orden $p + 1$ del violeta, en $0.8(p + 1)$: desde
$p = 2$ los órdenes se solapan. (c) $\delta = ax/D = \qty{2.5}{\micro m}$: faltan $\lambda = \delta/(p +
\tfrac12)$: $714$, $556$ y \qty{455}{nm}. (d) Centro blanco; luego franjas con un
borde interior azul y otro exterior rojo; hacia el tercer orden los colores
se mezclan en blanco.
\end{solution}

\begin{solution}{exo:b2:two-wave-interference:6}
(a) $b < \lambda L/4a = \qty{59}{\micro m}$. (b) $\pi ab/\lambda L = 6.7$: contraste $|\sin6.7|/6.7
= 0.06$. (c) Una fuente de un centímetro abarca cientos de interfranjas de
desplazamiento: sin franjas. (d) $a < \lambda/\theta = \qty{61}{\micro m}$.
\end{solution}

\begin{solution}{exo:b2:two-wave-interference:7}
(a) $a = 2h = \qty{0.4}{mm}$: $i = \lambda D/a = \qty{1.25}{mm}$. (b) En $x = 0$ los caminos
son iguales y sin embargo la franja es oscura: la reflexión añade $\pi$ (inversión
de signo con incidencia rasante sobre el medio más denso). (c) $\delta = ax/D <
\qty{50}{\micro m}$: $100$ franjas. (d) El mar es el espejo y la antena de lo alto del
acantilado, la pantalla: al subir la fuente, las franjas desfilan y
dan su elevación.
\end{solution}

\begin{solution}{exo:b2:two-wave-interference:8}
(a) $2d = 42\lambda$: $d = \qty{12.4}{\micro m}$. (b) $100/42 = \qty{2.4}{mm}$. (c) $2nd = p\lambda$:
$56$ franjas. (d) Las franjas se mueven hacia el extremo alejado y se separan a medida que $d$
disminuye; un cuarto de franja son $\lambda/8 = \qty{74}{nm}$ sobre $d$.
\end{solution}

\begin{solution}{exo:b2:two-wave-interference:9}
(a) Solo la primera (aire $\to$ aceite, más denso); aceite $\to$ agua va hacia un índice
menor: $\delta = 2ne + \lambda/2$, brillante en $2ne = (p + \tfrac12)\lambda$. (b) $2ne =
\qty{1160}{nm}$: reforzadas $773$ (borde del visible) y \qty{464}{nm} (azul);
canceladas $1160/p$: \qty{580}{nm} (amarillo): azul. (c) $2ne = \lambda/2$: $e = \qty{72}{nm}$.
(d) $\delta = 2ne\cos r$ decrece con el ángulo: las longitudes de onda reforzadas
se mueven hacia el azul. (e) Entonces ambas reflexiones invierten el signo, sin
$\lambda/2$ suplementaria: brillante en $2ne = p\lambda$ y oscuro en $(p + \tfrac12)\lambda$ --- la
\omterm{prop:b2:wave-interfaces:optics}{capa antirreflejante} de cuarto de onda.
\end{solution}

\begin{solution}{exo:b2:two-wave-interference:10}
(a) $(n - 1)e = \qty{10}{\micro m}$, $20$ franjas, hacia la rendija cubierta (donde
el camino geométrico debe acortarse \qty{10}{\micro m} para restaurar $\delta = 0$):
$x = (n - 1)eD/a = \qty{2}{cm}$. (b) La franja blanca se sitúa allí; la
dispersión del vidrio ($\Delta n \approx 0.01$ en el visible, es decir, \qty{0.2}{\micro m} de camino)
queda por debajo de $\ell_c$: sigue visible, algo coloreada. (c) Mídase el
desplazamiento de la franja blanca: $e = (\text{desplazamiento})\,a/(n - 1)D$. (d) $e(n - 1)\theta^2
(1 + 1/n)/2 = \qty{63}{nm}$ a \qty{5}{\degree}: un octavo de franja.
\end{solution}

\begin{solution}{exo:b2:two-wave-interference:11}
(a) $\theta = \qty{2.3e-7}{rad}$: $a = 1.22\lambda/\theta = \qty{3.1}{m}$. (b) Su límite $1.22\lambda/d =
\qty{0.057}{''}$ queda por encima del diámetro de la estrella, y la atmósfera difumina
hasta $1''$. (c) $1.22\lambda/100 = \qty{7e-9}{rad} = \qty{1.4}{mas}$. (d) Las franjas exigen
$\delta < \ell_c = \lambda^2/\Delta\lambda = \qty{32}{\micro m}$.
\end{solution}

\begin{solution}{exo:b2:two-wave-interference:12}
(a) $\delta = 2ne\cos r + \lambda/2$; $\cos r = \sqrt{1 - \sin^2i/n^2} \approx 1 - i^2/2n^2$.
(b) $2ne/\lambda = 10186.8$: el centro, con $p + \tfrac12$ entre $10186.5$ y
$10187.5$, no es ni brillante ni oscuro. (c) $i_p^2 = 2n^2[1 - (p + \tfrac12)\lambda/
2ne]$ para $p = 10186$, $10185$, $10184$: $i = 0.011$, $0.024$, \qty{0.032}{rad}, $r = fi
= 5.4$, $11.8$, \qty{15.8}{mm}. (d) El ángulo $i$ por sí solo fija el orden: cada
punto de la fuente alimenta el mismo anillo con el mismo ángulo, así que los anillos, vistos
en el infinito, sobreviven a una fuente amplia.
\end{solution}

\begin{solution}{pb:b2:two-wave-interference:1}
\textbf{1.} Cada espejo forma de $S$ una imagen por simetría a la distancia $d$ de
la línea de contacto; las dos imágenes están sobre la circunferencia de radio $d$,
separadas por el ángulo $2\alpha$: $a = 2\alpha d = \qty{2}{mm}$.

\textbf{2.} $S_1$ y $S_2$ son dos imágenes de la misma vibración: sus
saltos de fase son los mismos.

\textbf{3.} $i = \lambda(d + D)/a = 589 \times 10^{-9} \times 1.7/2 \times 10^{-3} = \qty{0.50}{mm}$: diez
franjas en \qty{5}{mm}.

\textbf{4.} $\delta = ax/(d + D) = \qty{2.9}{\micro m}$ en el borde, $p = 5$: muy por debajo de
$\ell_c/\lambda = 34\,000$.

\textbf{5.} Desplazamiento $bD'/d = 20 \times 10^{-6} \times 1.7/0.2 = \qty{0.17}{mm}$, un tercio de
$i$: contraste $|\operatorname{sinc}(\pi ab/\lambda d)| = \operatorname{sinc}(1.07) = 0.82$.

\textbf{6.} $\pi ab/\lambda d = 5.3$: $|\sin5.3|/5.3 = 0.16$.

\textbf{7.} Amplitudes $\sqrt{0.9}$, $\sqrt{0.8}$: $V = 2\sqrt{0.72}/1.7 = 0.998$.

\textbf{8.} $(n - 1)e = \qty{5}{\micro m}$: $8.5$ franjas; se sigue sin dificultad con
luz de sodio; con luz blanca la franja blanca se mueve $8.5i = \qty{4.2}{mm}$
y sigue ahí (la dispersión del vidrio cuesta una fracción de
micrómetro).

\textbf{9.} Una franja central blanca, dos o tres franjas coloreadas a
cada lado y luego un blanco uniforme.

\textbf{10.} Aire $\to$ aceite invierte, aceite $\to$ agua (índice menor) no:
$\delta = 2ne + \lambda/2$; brillante para $2ne = (p + \tfrac12)\lambda$.

\textbf{11.} $e = \lambda/4n$: rojo \qty{112}{nm}, verde \qty{95}{nm}, violeta \qty{72}{nm}.

\textbf{12.} $2ne = \qty{1740}{nm}$: reforzadas $1740/(p + \tfrac12)$: $696$ (rojo), $497$
(azul verdoso), \qty{387}{nm}; canceladas $1740/p$: $580$ (amarillo), \qty{435}{nm}:
una mezcla púrpura de rojo y azul verdoso.

\textbf{13.} El espesor varía a lo largo de la carretera; cada color es brillante
allí donde $2ne$ toma el valor adecuado, así que las bandas son curvas de igual
espesor.

\textbf{14.} Los órdenes bajan: los colores recorren la secuencia
hacia el extremo delgado; cuando $e \to 0$, la única inversión de signo deja
$\delta = \lambda/2$ para todos los colores: oscuro --- como para una película de jabón en el aire.

\textbf{15.} Unos pocos nanómetros: $2ne \ll \lambda$ para todos los colores, las dos
reflexiones están en antifase y se cancelan: negro.

\textbf{16.} $\sin r = \sin60^\circ/1.45$, $r = \qty{37}{\degree}$, $\cos r = 0.80$: $2ne\cos r
= \qty{1390}{nm}$, brillante en $928$, $557$ y \qty{398}{nm}: el rojo de la pregunta 12 se ha
vuelto verde amarillento --- hacia el azul.

\textbf{17.} Cuando $2ne$ supera $\ell_c \approx \qty{1}{\micro m}$, es decir, más allá de unos
\qty{0.4}{\micro m} de aceite, los órdenes se solapan y los colores se apagan en gris.

\textbf{18.} Los dos rayos reflejados por la lámina para cualquier dirección
de la fuente se reencuentran en la lámina: las franjas están localizadas allí y
una fuente amplia y multidireccional se limita a añadir patrones idénticos.

\textbf{19.} $\ell_c/\lambda$: $34\,000$ con sodio, $1.7 \times 10^7$ con el láser.

\textbf{20.} $b < \lambda d/4a$: \qty{15}{\micro m} para los espejos y \qty{59}{\micro m} para
las rendijas.

\textbf{21.} Dibújense dos puntos de la fuente: para cada uno, los dos rayos reflejados
se recombinan en el mismo punto de la lámina, con la misma $\delta = 2ne\cos r$
con incidencia casi normal; el patrón de la lámina no se mueve, y el de Young sí.

\textbf{22.} $\lambda/\Delta\lambda = 589/0.02 \approx 30\,000$.

\textbf{23.} $550/30 \approx 18$ franjas.

\textbf{24.} La luz coherente dispersada por los granos de la pantalla interfiere
con fases al azar en el ojo: moteado.

\textbf{25.} Espejos: \omterm{prop:b2:two-wave-interference:young}{división del frente de onda}, franjas en todas partes
(no localizadas), número limitado por $\ell_c/\lambda$, y una fuente ancha las
borra. Lámina: \omterm{prop:b2:two-wave-interference:film}{división de amplitud}, franjas en la lámina, colores
limitados por $\ell_c \sim \qty{1}{\micro m}$, y una fuente ancha es inocua.
\end{solution}
